\noindent\hfill
\begin{minipage}[t]{0.6\textwidth}
    %RESUMEN DE LA TESIS PARA OPTAR AL \\
    %GRADO DE DOCTOR[A] EN INGENIERÍA \\
    %ELÉCTRICA, \\
    RESUMEN DE LA TESIS PARA OPTAR AL \\
    GRADO DE MAGÍSTER EN CIENCIAS DE LA \\ INGENIERÍA, MENCIÓN ELÉCTRICA, \\
    MEMORIA PARA OPTAR AL TÍTULO DE \\ INGENIER[O/A] CIVIL ELÉCTRIC[O/A] \\
    {POR: NOMBRE APELLIDO1 APELLIDO2} \\
    {FECHA: \the\year} \\
    {PROF. GUÍA: NOMBRE APELLIDO1 APELLIDO2} \\
    %{PROF. CO-GUÍA: NOMBRE COMPLETO CO-GUÍA}
\end{minipage}

\begin{center}
    TÍTULO DE LA TESIS
\end{center}

\noindent En esta sección se debe incluir un resumen de máximo doscientas (200) palabras de la tesis en un solo párrafo. Es posible incluir adicionalmente un resumen extendido en la sección de anexos si es que fuera necesario. En general, para construir esta sección se recomienda incluir una breve motivación que indique por qué el tópico es relevante; una definición clara y concisa del problema; una breve descripción del método para resolver el problema; uno o unos pocos resultados relevantes, incluyendo idealmente una métrica numérica verificable (por ejemplo, ``la eficiencia energética aumentó un 19\%''); y finalmente una conclusión que indique claramente los hallazgos de la investigación y el impacto de estos en el área. Para cumplir con el límite de palabras, se recomienda que cada uno de los elementos mencionados no supere las dos oraciones. Dado que al momento de leer el manuscrito la tesis ya está concluida, no se debe redactar en tiempo futuro (por ejemplo, ``en este trabajo se analizará''). Esta sección debe ser autocontenida; por ello, es importante asegurar la cohesión entre oraciones y no incluir referencias bibliográficas. Tampoco se deben incluir tablas ni figuras. En caso de utilizar una cantidad considerable de símbolos o acrónimos, se recomienda incluir una lista de nomenclatura al inicio del documento.

%%%%%%%%%%%%%%%%%%%%%%%%%%%%%%%%%%%%%%%
% SÓLO PARA TESIS EN INGLÉS 
% BORRAR EN CASO DE ESCRITURA EN ESPAÑOL
\newpage

\noindent\hfill
\begin{minipage}[t]{0.6\textwidth}
    %RESUMEN DE LA TESIS PARA OPTAR AL \\
    %GRADO DE DOCTOR[A] EN INGENIERÍA \\
    %ELÉCTRICA, \\
    RESUMEN DE LA TESIS PARA OPTAR AL \\
    GRADO DE MAGÍSTER EN CIENCIAS DE LA \\ INGENIERÍA, MENCIÓN ELÉCTRICA, \\
    MEMORIA PARA OPTAR AL TÍTULO DE \\ INGENIER[O/A] CIVIL ELÉCTRIC[O/A] \\
    {POR: NOMBRE APELLIDO1 APELLIDO2} \\
    {FECHA: \the\year} \\
    {PROF. GUÍA: NOMBRE APELLIDO1 APELLIDO2} \\
    %{PROF. CO-GUÍA: NOMBRE COMPLETO CO-GUÍA}
\end{minipage}

\begin{center}
    TÍTULO DE LA TESIS (INGLÉS)
\end{center}

\noindent En el caso de redactar la tesis en inglés, se debe incluir un resumen adicional en dicho idioma con su título en inglés, ya que el resumen en español (con título en español) es obligatorio independientemente del idioma de la tesis. Ambos resúmenes deben reflejar el mismo contenido. Sin embargo, dado que cada uno debe cumplir independientemente con el límite de doscientas (200) palabras, no es necesario que sean una traducción textual el uno del otro. Las mismas recomendaciones estructurales indicadas en el resumen anterior aplican para ambos.
